\section{Nota complementaria sobre grado y eigenvector}

Las actividades previas sugieren una distinción útil entre grado y centralidad de eigenvector. Un nodo puede tener muchos vecinos y, sin embargo, estar conectado principalmente a vértices periféricos; en ese caso, su grado será alto pero su eigenvector no necesariamente. Recíprocamente, un nodo con pocas conexiones directas puede obtener una centralidad espectral relativamente alta si sus vecinos están muy bien posicionados dentro del núcleo de la red.

En la minired del S\&P 500, esta observación ayuda a entender por qué algunos nodos con varios enlaces no destacan en eigenvector: sus conexiones no apuntan hacia la parte más influyente del subgrafo. Del mismo modo, explica por qué ciertos compradores centrales reciben puntuaciones espectrales altas aun cuando no concentran la mayor cantidad de enlaces salientes.

En la red de \textit{Star Wars Episode II}, la diferencia entre grado y eigenvector es menor porque el núcleo narrativo está más concentrado y las conexiones de los protagonistas recaen, en buena medida, sobre otros personajes también centrales. Aun así, la comparación sigue siendo útil para distinguir entre popularidad local y cercanía estructural al corazón del grafo.
